% ============================================================
% v2 — Capitolo 4: Il progetto — risultati e discussione
% Blocco II di IV. Budget: ~10 pagine.
% Ogni numero è ricalcolato da episodes.jsonl su disco
% (docs/thesis/data/campaign_consolidation.json). Le tabelle eval
% MLP e GNN della v1 sono fuse in un'unica tabella (tab:eval).
% ============================================================
\chapter{Il progetto --- risultati e discussione}
\label{ch:risultati}

Questo capitolo riporta i risultati della campagna lungo i quattro assi
di confronto, aggiunge l'asse architetturale, chiude con i verdetti
sulle tre ipotesi e ne discute il significato. Tutti i numeri sono
ricalcolati dai record a livello di agente su disco secondo le regole
della \cref{sec:metriche}.

% ------------------------------------------------------------
\section{Integrità della campagna}
\label{sec:integrita}

Le quattro run sono state lanciate sullo stesso trunk con seed 999, e i
diff di configurazione confermano che le celle differiscono solo nei
flag di regime ed encoder del critic. Tutti e otto i log --- quattro di
addestramento e quattro di valutazione --- superano ogni condizione di
integrità della \cref{sec:metriche}: esattamente sei record per
episodio, zero righe malformate, zero valori non finiti. Le run hanno
prodotto tra 3{.}168 e 3{.}293 episodi ciascuna, poco più di 19{.}000
record; l'unica anomalia a registro è benigna, poiché la valutazione
GNN-curriculum ha mantenuto 401 episodi dopo la deduplica, uno
spostamento di peso dello $0.25$\% senza effetto alla precisione
riportata. L'esposizione di livello realizzata rispecchia infine i due
design: entrambe le run batch restano entro $\approx 2$ punti
percentuali dall'uniforme $1/3$, con tutti e tre i livelli presenti
entro i primi venti episodi, mentre le run curriculum concentrano
l'esposizione su easy e medium all'inizio e raggiungono hard solo dopo i
rispettivi sblocchi, chiudendo a 32.5/45.8/21.6\% (MLP) e
32.4/37.4/30.1\% (GNN).

% ------------------------------------------------------------
\section{Efficienza campionaria}
\label{sec:efficienza}

\begin{figure}[htbp]
  \centering
  \includegraphics[width=\textwidth]{learning_curves}
  \caption{Tasso di successo a finestra (100 episodi) rispetto ai passi
  d'ambiente, tutte e quattro le celle.}
  \label{fig:curve}
\end{figure}

Una avvertenza va premessa alla \cref{fig:curve}: ogni curva è misurata
sul \emph{mix di livelli concorrente} della propria run --- gated per le
celle curriculum, uniforme per batch --- quindi i confronti verticali
tra bracci non sono omogenei finché il curriculum non ha sbloccato tutti
i livelli.

\paragraph{Fase iniziale: decollo indistinguibile.}
Tutte e quattro le celle salgono alla stessa velocità. L'SR veicoli a
finestra di 100 episodi supera per la prima volta il 30\% dopo
137--142K passi e il 50\% dopo 162--175K passi \emph{in ogni cella} ---
uno scarto sotto l'8\%, ben dentro il rumore. La parte notevole è
\emph{su che cosa} ciascuna policy ha superato quelle soglie: le celle
curriculum su esposizione solo-easy, le celle batch sul mix uniforme
completo, percorsi da \SI{100}{m} inclusi. A parità di budget, il
decollo iniziale non è stato il fattore discriminante tra i regimi.

\paragraph{Tempi di sblocco.}
Medium si è sbloccato praticamente nello stesso punto in entrambe le run
curriculum (al 26.1\% degli episodi in entrambe). Hard, invece, si è
sbloccato molto prima sotto il critic GNN: al 58.0\% dell'addestramento
contro il 71.0\% per MLP --- quest'ultimo coerente con lo scatto del cap
di pressione al 70\% del budget più che con il gate di competenza,
mentre la cella GNN ha raggiunto il gate hard \emph{per competenza}, e
ha quindi visto $\approx 40$\% episodi hard in più. Da $\approx 0.5$M
passi, infine, le curve batch si assestano sotto le due curve
curriculum, che si seguono da vicino; tutte derivano poi verso il basso
a fine addestramento (\cref{sec:dinamica}).

% ------------------------------------------------------------
\section{Prestazione finale}
\label{sec:prestazione}

\begin{table}[htbp]
  \centering\small
  \caption{I due strati di misura. Sopra: metriche cumulative di
  addestramento. Sotto: valutazione finale del checkpoint, 400 episodi
  per cella (401 per GNN-curriculum, \cref{sec:integrita}), l'unico
  confronto a distribuzione appaiata tra i bracci. Valori a livello di
  agente in \%; velocità in km/h (veicoli) e m/s (pedoni).}
  \label{tab:cumulative}
  \begin{tabular}{@{}lrrrrrrrr@{}}
    \toprule
    & \multicolumn{5}{c}{Veicoli} & \multicolumn{3}{c}{Pedoni} \\
    \cmidrule(lr){2-6}\cmidrule(l){7-9}
    Cella & SR & s$+$t & coll. & off. & vel. & SR & s$+$t & vel. \\
    \midrule
    MLP-curr.\ & 57.28 & 30.48 & 8.84 & 3.41 & 10.2 & 56.59 & 42.56 & 1.96 \\
    MLP-batch  & 46.31 & 23.61 & 17.33 & 12.76 & 18.7 & 57.06 & 42.03 & 1.61 \\
    GNN-curr.\ & 56.55 & 20.84 & 13.37 & 9.23 & 18.0 & 60.50 & 38.58 & 1.91 \\
    GNN-batch  & 49.41 & 22.75 & 17.56 & 10.28 & 26.6 & 58.03 & 41.11 & 1.45 \\
    \bottomrule
  \end{tabular}

  \vspace{1.2em}

  \begin{tabular}{@{}lrrrrrrrrr@{}}
    \toprule
    & \multicolumn{5}{c}{SR veicoli per scenario} &
      \multicolumn{3}{c}{Veicoli, compl.} & Ped. \\
    \cmidrule(lr){2-6}\cmidrule(lr){7-9}\cmidrule(l){10-10}
    Cella & easy & med. & hard & test & \textbf{tutti} & s$+$t & coll. & off. & SR \\
    \midrule
    MLP-curr.\ & 60.33 & 51.67 & 44.33 & 37.00 & \textbf{48.33} & 38.25 & 10.42 & 3.00 & 47.00 \\
    MLP-batch  & 41.33 & 31.00 & 29.00 & 18.00 & \textbf{29.83} & 31.17 & 26.00 & 13.00 & 48.92 \\
    GNN-curr.\ & 52.48 & 42.33 & 36.00 & 35.67 & \textbf{41.65} & 25.10 & 22.94 & 10.31 & 59.02 \\
    GNN-batch  & 55.67 & 42.00 & 33.33 & 25.67 & \textbf{39.17} & 23.83 & 26.92 & 10.08 & 48.67 \\
    \bottomrule
  \end{tabular}
  \\[0.3em]
  {\footnotesize Scenari easy/medium/hard su Town03, \emph{test} su Town05.}
\end{table}

\paragraph{Il contrasto primario, coppia MLP.}
Nella valutazione finale a distribuzione appaiata
(\cref{tab:cumulative}, pannello inferiore) la policy veicolo addestrata
a curriculum completa il 48.33\% dei percorsi contro il 29.83\% del
batch: $+18.5$ punti percentuali, con lo stesso segno in \emph{tutti e
quattro} gli scenari ($+19.0$, $+20.7$, $+15.3$, $+19.0$\,pp) e nel
cumulativo di addestramento ($+11.0$\,pp). Il margine è circa da due a
quattro volte la banda delle repliche a seed identico della
\cref{sec:gate}, e la sua coerenza tra scenari, livelli di difficoltà e
strati di misura è ciò che il design appaiato può offrire al posto della
replicazione dei seed.

\paragraph{Due modi diversi di fallire.}
I regimi producono guide riconoscibilmente diverse
(\cref{fig:fallimenti}). La policy batch è quasi due volte più veloce
(18.7 contro 10.2 km/h in addestramento) ma converte quella mobilità in
contatto: collisione 17.3\% contro 8.8\% e offroad 12.8\% contro 3.4\%
in addestramento, che si allargano a 26.0\% contro 10.4\% di collisione
in valutazione. La policy curriculum fallisce invece principalmente per
stallo (stuck$+$timeout 38.3\% contro 31.2\% in valutazione). Nessun
profilo domina l'altro su ogni asse --- il curriculum scambia velocità
per sicurezza --- ma sul criterio di successo rigoroso lo scambio si
risolve chiaramente a suo favore. I pedoni, per contro, sono
\emph{piatti} tra i bracci ovunque, entro $\pm 6$\,pp in ogni scenario:
il regime, in questa piattaforma, è una variabile della policy
\emph{veicolo}, e la struttura dei fallimenti pedonali è ripresa nella
\cref{sec:pedoni}.

\begin{figure}[t]
  \centering
  \includegraphics[width=\textwidth]{failure_modes}
  \caption{Composizione delle terminazioni dei veicoli per cella, in
  addestramento (cumulativo) e alla valutazione finale. Le due
  grammatiche del fallimento sono visibili a colpo d'occhio: dominata
  dallo stallo in MLP-curriculum, dominata dal contatto nelle celle
  batch, con le celle GNN in posizione intermedia.}
  \label{fig:fallimenti}
\end{figure}

\paragraph{Le celle GNN a colpo d'occhio.}
In addestramento la cella GNN-curriculum atterra entro il rumore di
MLP-curriculum sull'SR principale (56.55 contro 57.28\%) ma
ridistribuisce la massa dei fallimenti: meno stallo ($-9.6$\,pp), più
contatto e velocità molto più alta (18.0 contro 10.2 km/h) --- più
vicina alla mobilità delle celle batch, con il successo di una cella
curriculum. GNN-batch spinge le stesse tendenze oltre (26.6 km/h, la
cella più veloce) \emph{battendo} al contempo MLP-batch di $+3.1$\,pp in
addestramento e $+9.3$\,pp in valutazione: l'encoder del critic plasma
il profilo di rischio e interagisce con il regime
(\cref{sec:asse-arch}).

% ------------------------------------------------------------
\section{Dinamica entro la run}
\label{sec:dinamica}

Ogni curva della \cref{fig:curve} deriva verso il basso dopo il picco
iniziale. Per le celle curriculum questo è in parte meccanico --- il mix
concorrente si indurisce a ogni sblocco --- quindi la
\cref{tab:quartili} separa i due effetti: il pannello superiore riporta
le traiettorie per quartile complessive, quello inferiore la traiettoria
\emph{sui soli episodi easy}, dove il task è costante per costruzione.

\begin{table}[t]
  \centering\small
  \caption{SR veicoli per quartile di addestramento (\%). Sopra: tutti i
  livelli (confuso dal mix per le celle curriculum). Sotto: solo episodi
  easy, controllato per livello (easy resta $\approx 30$\% del mix dopo
  gli sblocchi). $\Delta$ = Q4 $-$ Q1.}
  \label{tab:quartili}
  \begin{tabular}{@{}lrrrrr@{}}
    \toprule
    Cella & Q1 & Q2 & Q3 & Q4 & $\Delta$ \\
    \midrule
    \multicolumn{6}{@{}l}{\emph{Tutti i livelli}} \\
    MLP-curr.\ & 68.9 & 61.1 & 54.4 & 44.7 & $-24.2$ \\
    MLP-batch  & 58.7 & 45.9 & 41.8 & 38.7 & $-20.0$ \\
    GNN-curr.\ & 66.6 & 62.2 & 54.5 & 42.9 & $-23.7$ \\
    GNN-batch  & 50.5 & 54.6 & 47.5 & 45.1 & $-5.4$ \\
    \midrule
    \multicolumn{6}{@{}l}{\emph{Solo livello easy}} \\
    MLP-curr.\ & 68.9 & 69.3 & 59.6 & 52.4 & $-16.5$ \\
    MLP-batch  & 64.2 & 52.0 & 48.9 & 45.6 & $-18.6$ \\
    GNN-curr.\ & 66.6 & 67.9 & 56.1 & 55.8 & $-10.8$ \\
    GNN-batch  & 59.7 & 62.8 & 54.3 & 53.3 & $-6.4$ \\
    \bottomrule
  \end{tabular}
\end{table}

\paragraph{Erosione, non solo indurimento del mix.}
Il pannello controllato per livello mostra che l'indurimento del mix
\emph{non} è tutta la storia: la prestazione a livello easy si erode
entro ogni run, da $-6.4$\,pp (GNN-batch, la cella più stabile) a
$-18.6$\,pp (MLP-batch). Le celle batch rendono il caso in modo pulito,
perché la loro distribuzione dei task è costante e il loro declino è
quindi interamente degradazione entro la run; i pedoni si erodono ancora
di più, fino a $-31.8$\,pp. Due meccanismi candidati sono compatibili
con i dati e non separabili a $n{=}1$: la co-adattazione tra classi,
poiché il comportamento dei pedoni degrada in parallelo e alimenta il
canale di hazard dei veicoli, e l'annealing dell'entropia --- che però
si completa solo dentro Q4, mentre l'erosione inizia già in Q3, e non
può quindi esserne l'unica causa. La conclusione pulita è
architetturale: sotto il critic GNN l'erosione è dimezzata o meglio,
qualunque ne sia il driver.

\paragraph{Due grammatiche del fallimento, un attrattore comune.}
La velocità media ordina le quattro celle --- 10.2 (MLP-curr.), 18.0
(GNN-curr.), 18.7 (MLP-batch), 26.6 km/h (GNN-batch) --- e i fallimenti
da contatto seguono lo stesso ordine, mentre lo stallo anti-correla. Ciò
che \emph{non} varia è la struttura interna dello stallo: in tutte e
quattro le celle le terminazioni stuck si decompongono quasi
identicamente e i passi medi di non-progresso coprono una banda stretta.
Regimi e critic cambiano \emph{quanto} una policy stalla, non
\emph{come} stalla. Coerentemente, il profilo lento-e-in-stallo
appartiene a esattamente una cella, MLP-curriculum: la sua gemella sotto
il critic GNN è quasi due volte più veloce con un terzo di stallo in
meno, allo stesso SR principale. L'equilibrio cauto documentato durante
lo sviluppo (\cref{sec:gate}) non è dunque una proprietà
dell'addestramento a curriculum in quanto tale, ma un bacino in cui la
run con critic MLP si è assestata e quella con critic GNN no: il regime
da solo non determina l'equilibrio.

% ------------------------------------------------------------
\section{Generalizzazione: da Town03 a Town05}
\label{sec:generalizzazione}

Lo scenario test colloca le stesse policy su Town05 --- layout a griglia
quadrata e più corsie, strutturalmente diverso da Town03 --- per 100
episodi con target di \SI{80}{m} e traffico come in addestramento.

\begin{figure}[t]
  \centering
  \includegraphics[width=0.70\textwidth]{town05_bars}
  \caption{SR veicoli: media degli scenari Town03 contro test Town05,
  per cella. Le celle curriculum calano meno; la coppia GNN è
  indistinguibile in-domain (43.6 contro 43.7\%) e si separa di 10\,pp
  sulla mappa mai vista.}
  \label{fig:town05}
\end{figure}

\paragraph{L'ordinamento sopravvive al cambio di mappa.}
Il curriculum completa il 37.0\% dei percorsi su Town05 contro il 18.0\%
del batch ($+19.0$\,pp): essenzialmente lo stesso margine in-domain. Il
calo assoluto dalla media degli scenari Town03 è quasi identico per i
due bracci ($-15.1$ contro $-15.8$\,pp), quindi il vantaggio del
curriculum non è un'idiosincrasia di Town03: si trasferisce per intero
su strade mai viste. Letta contro H2, la copertura iniziale più ampia
del regime batch non gli ha comprato \emph{alcun} vantaggio di
generalizzazione --- semmai il calo relativo è peggiore, dalla sua base
più bassa ($-47$\% contro $-29$\%). Ogni braccio, inoltre, fallisce su
Town05 nel modo in cui fallisce a casa, solo di più: la policy
curriculum stalla (stuck$+$timeout 50.7\%), la policy batch si schianta
(collisione 37.3\%).

\paragraph{Sotto il critic GNN il vantaggio esiste \emph{solo} fuori distribuzione.}
La coppia GNN affila il punto. In-domain i due bracci sono gemelli
statistici --- SR veicoli medio sugli scenari Town03 43.6\% (curriculum)
contro 43.7\% (batch) --- eppure su Town05 si separano di $+10.0$\,pp
(35.7 contro 25.7\%), con il calo più piccolo di ogni cella per
GNN-curriculum ($-7.9$\,pp) contro $-18.0$\,pp per GNN-batch. Dove la
coppia MLP mostrava un vantaggio del curriculum che \emph{si
trasferisce}, la coppia GNN ne mostra uno che \emph{compare} solo nel
trasferimento: evidenza con lo stesso segno, veicolata da due firme
diverse. I pedoni, infine, si trasferiscono con facilità --- $\approx
30$\,pp sopra i risultati hard di entrambi i bracci su Town03 ---
indicando ancora la fattibilità dei percorsi, non la generalizzazione
della policy, come loro collo di bottiglia (\cref{sec:pedoni}).

% ------------------------------------------------------------
\section{Asse architetturale: critic MLP contro GNN}
\label{sec:asse-arch}

L'asse secondario sostituisce solo l'encoder del critic; actor, reward,
ambiente e budget sono identici, quindi le differenze tra celle MLP e
GNN misurano come la \emph{stima del valore} plasma ciò che gli stessi
actor apprendono.

\begin{table}[t]
  \centering\small
  \caption{L'interazione architettura$\times$regime: SR veicoli
  curriculum-meno-batch (pp) per strato di misura, entro ciascuna
  architettura. Il delta di valutazione dei pedoni è aggiunto per
  completezza.}
  \label{tab:interazione}
  \begin{tabular}{@{}lrr@{}}
    \toprule
    $\Delta$(curr.\ $-$ batch) & Coppia MLP & Coppia GNN \\
    \midrule
    Cumulativo di addestramento (veic.) & $+10.97$ & $+7.14$ \\
    Valutazione complessiva (veic.) & $+18.50$ & $+2.48$ \\
    Valutazione, media Town03 (veic.) & $+18.33$ & $-0.06$ \\
    Valutazione, Town05 (veic.) & $+19.00$ & $+10.00$ \\
    Valutazione complessiva (ped.) & $-1.92$ & $+10.35$ \\
    \bottomrule
  \end{tabular}
\end{table}

\paragraph{L'effetto di regime dipende dal critic.}
La \cref{tab:interazione} condensa il risultato centrale della campagna.
Sotto il critic MLP il curriculum batte il batch con un margine ampio e
coerente tra gli strati; sotto il critic GNN lo stesso contrasto quasi
svanisce, con $+2.5$\,pp complessivi in valutazione --- dentro la banda
di rumore a seed identico --- ed esattamente zero in-domain. L'unica
componente del vantaggio del curriculum che \emph{si replica in entrambe
le architetture} è quella out-of-distribution: $+19.0$\,pp (MLP) e
$+10.0$\,pp (GNN) su Town05. H3 è quindi supportata nella sua forma
forte: l'encoder del critic cambia non solo l'entità, ma il \emph{luogo}
dell'effetto di regime.

\paragraph{Da dove viene l'interazione.}
La compressione è prodotta da due effetti architetturali asimmetrici,
uno per regime. Entro il batch il critic GNN \emph{aiuta}: $+3.1$\,pp in
addestramento e $+9.3$\,pp in valutazione rispetto a MLP-batch, con la
dinamica entro la run più stabile della campagna. È coerente con il
meccanismo che motiva i critic CTDE in primo luogo: sotto una
distribuzione mista con sei agenti interagenti, la stima strutturata del
valore dovrebbe aiutare di più dove il segnale è più rumoroso, cioè il
regime batch dal passo~0. Entro il curriculum, invece, il critic GNN
\emph{rimodella} più che migliorare: evita l'equilibrio difensivo,
sblocca hard per competenza anziché per cap di pressione, vede $\approx
40$\% episodi hard in più --- e converte tutto questo in mobilità e
contatto anziché in successo (SR $-6.7$\,pp rispetto a MLP-curriculum).
I pedoni mostrano l'interazione speculare: il critic GNN li solleva
sotto curriculum ($+12.0$\,pp in valutazione, la miglior cella pedoni
della campagna al 59.0\%) ma non sotto batch.

Il pattern è internamente coerente su quattro strati di misura ed
entrambe le classi di agenti, ma le sue grandezze portano la consueta
qualifica $n{=}1$ (\cref{sec:minacce}): il delta complessivo di
$+2.5$\,pp della coppia GNN esplicitamente \emph{non} è rivendicato come
vantaggio del curriculum, e il risultato rivendicato è il suo contrasto
con la coppia MLP più la componente Town05 replicata tra le
architetture.

% ------------------------------------------------------------
\section{Limiti strutturali dei pedoni}
\label{sec:pedoni}

I risultati dei pedoni vanno letti contro un vincolo che nessun regime
di addestramento può rimuovere: la topologia dei marciapiedi di Town03
spesso non può fornire il task nominale. Le diagnostiche di generazione
dei percorsi pedonali sono infatti statisticamente indistinguibili tra
le quattro celle --- percorsi under-target 55.1--55.9\% dei record, flag
too-short 7.9--9.2\%, percorsi da respawn 13.4--14.1\% --- e questa
\emph{invarianza} è il punto: sono proprietà della mappa e del planner,
non delle policy, e fissano quindi un soffitto comune a ogni cella.

L'SR pedoni cala poi monotonicamente con il livello in ogni cella (in
addestramento easy 70--83\%, medium 53--60\%, hard 29--44\%) e la
valutazione a distribuzione appaiata riproduce il gradiente. Un collasso
che sopravvive a due architetture, due regimi e due strati di misura
mentre le diagnostiche restano costanti è guidato dalla geometria del
task, e il controfattuale Town05 rende concreta l'attribuzione: con
target di \SI{80}{m} l'SR pedoni balza a 54--71\% in ogni cella, da 15 a
31\,pp \emph{sopra} il risultato hard su Town03 e per giunta su una
mappa mai vista --- se il collasso fosse un limite di policy, strade
nuove dovrebbero nuocere, non aiutare. Entro l'inviluppo di fattibilità
le differenze di policy restano comunque visibili, e una cella si separa
dal gruppo in valutazione (GNN-curriculum, 59.0\% contro 47--49\%
altrove). La conseguenza metodologica resta però vincolante: i tassi di
successo medium e hard dei pedoni non sono una misura pulita della
qualità della policy, quindi la tesi confronta i pedoni solo entro
livello e tra bracci.

% ------------------------------------------------------------
\section{Verdetti sulle ipotesi}
\label{sec:verdetti}

\begin{table}[t]
  \centering\small
  \caption{Ipotesi contro evidenza (contrasti appaiati, $n{=}1$ per
  cella; delta in pp di SR veicoli salvo diversa indicazione).}
  \label{tab:ipotesi}
  \begin{tabular}{@{}p{2.0cm}p{4.3cm}p{2.3cm}p{4.2cm}@{}}
    \toprule
    Ipotesi & Predizione & Verdetto & Evidenza \\
    \midrule
    H1 (curriculum) & migliore efficienza campionaria \emph{iniziale} & rigettata &
      decollo identico nelle quattro celle (50\% a 162--175K passi) \\
    & prestazione finale più alta & supportata sotto MLP, non sotto GNN
      in-domain & $+18.5$ (MLP) contro $+2.5$ (GNN) in valutazione \\
    & rischio di plateau od oblio su hard & rimodellato &
      erosione tardiva reale ma \emph{di piattaforma}: $\Delta$Q1$\to$Q4
      per livello da $-6$ a $-19$\,pp \\
    \addlinespace
    H2 (batch) & piena copertura dal passo 0 & confermata (meccanica) &
      tutti i livelli entro i primi 20 episodi \\
    & segnale iniziale più rumoroso & non visibile &
      stesso decollo, sul mix più difficile \\
    & copertura $\to$ robustezza su mappa nuova & rigettata &
      calo su Town05 uguale o peggiore in entrambe le coppie
      ($-15.8$/$-18.0$ contro $-15.1$/$-7.9$\,pp) \\
    \addlinespace
    H3 (interazione) & l'effetto di regime può dipendere dall'encoder
      del critic & supportata (forma forte) &
      effetto in-domain $+18.3$ (MLP) contro $-0.1$ (GNN); solo la
      componente Town05 si replica \\
    \bottomrule
  \end{tabular}
\end{table}

Due cose sopravvivono a ogni cella, ogni strato ed entrambe le
architetture: il vantaggio del curriculum sulla mappa mai vista, e il
ruolo del regime come variabile della policy \emph{veicolo} che rimodella
la composizione dei fallimenti. Ciò che non è sopravvissuto è
l'aspettativa popolare da entrambi i lati: il curriculum non ha appreso
più in fretta all'inizio, e la copertura più ampia del batch non ha
comprato robustezza.

% ------------------------------------------------------------
\section{Discussione}
\label{sec:discussione}

\paragraph{Non la velocità di apprendimento, ma la selezione dell'equilibrio.}
L'asse su cui i curricula sono classicamente attesi ripagare non ha
mostrato nulla. Una lettura plausibile è che il design della reward di
questa piattaforma risolva già il problema per cui i curricula vengono
di solito ingaggiati: con bonus densi per waypoint e shaping sulla
distanza (\cref{sec:reward}), anche i percorsi da \SI{100}{m}
forniscono segnale di apprendimento fin dai primi episodi, quindi i task
facili non servono per l'avvio. È una \emph{condizione di scope} su H1,
più che una confutazione del curriculum learning in generale. Ciò che
l'ordinamento ha cambiato è \emph{quale policy lo stesso apprendista è
diventato}, e la meccanica è visibile nei vincoli: completare un
percorso da \SI{100}{m} entro l'orizzonte di mille passi richiede una
velocità media sana, quindi una policy esposta ai percorsi hard dal
passo~0 deve guidare veloce per riuscire del tutto --- e paga in
collisioni --- mentre una policy che spende il suo primo quarto di
addestramento su percorsi da \SI{30}{m} può riuscire lentamente, e
cementa la cautela. L'ordinamento agisce dunque da \emph{selettore di
equilibri} sul trade-off velocità--sicurezza più che da acceleratore, e
le celle GNN sostengono la lettura dall'altro lato: con uno stimatore di
valore diverso lo stesso curriculum ha selezionato un equilibrio più
veloce. Va infine osservato che l'erosione entro la run è di piattaforma
e che l'ordinamento non la causa né la previene: i numeri cumulativi
\emph{finali} sottostimano i picchi di metà addestramento in ogni
configurazione, e la selezione del checkpoint cambierebbe le
conclusioni.

\paragraph{Generalizzazione: il beneficio portabile.}
L'unico vantaggio replicato in entrambe le architetture è il margine su
Town05, ottenuto \emph{nonostante} la copertura in-domain più ampia del
batch. Un meccanismo provvisorio è coerente con entrambe le coppie: la
policy batch si co-adatta per l'intero budget a una singola
distribuzione stazionaria, mentre la policy curriculum deve restare
competente attraverso una sequenza \emph{mutevole} di distribuzioni ---
una forma blanda di regolarizzazione implicita per non-stazionarietà. A
$n{=}1$ per cella questa resta un'ipotesi; ciò che i dati stabiliscono è
la direzione e la sua replicazione tra architetture.

\paragraph{La copertura parallela dei task rivisitata.}
La cornice a mondo condiviso è stata adottata per il throughput, con
campioni correlati e non-stazionarietà reciproca accettati come costi;
col senno di poi la campagna mostra quei ``costi'' fare lavoro
scientifico di prim'ordine. L'equilibrio difensivo esiste \emph{solo}
perché le reward dei veicoli sono condizionate a un canale di hazard
alimentato dal comportamento dei pedoni, e il risultato architetturale
punta nella stessa direzione: sostituire un critic piatto con uno che
modella gli agenti come nodi interagenti ha cambiato gli esiti in ogni
cella, cosa impossibile se la struttura tra agenti fosse trascurabile.
In assenza di un controllo a un agente per mondo, però, il beneficio di
throughput resta un argomento di design e non una misura.

% ------------------------------------------------------------
\section{Minacce alla validità e limiti di scope}
\label{sec:minacce}

\paragraph{Non-determinismo del simulatore e dimensione campionaria.}
La minaccia dominante è quantificata anziché assunta: tre run a
lunghezza piena con codice, configurazione e seed byte-identici hanno
prodotto tassi di successo cumulativi dei veicoli tra 52.2 e 62.4\%
(\cref{sec:gate}). Ogni effetto è letto contro questa banda di
$\sim$10\,pp, e le differenze al suo interno non sono mai attribuite a
un fattore sperimentale. A ciò si somma il fatto che ogni cella contiene
una sola run (\cref{sec:campagna}): l'appaiamento dei seed rende i
contrasti internamente puliti, e il contrasto primario porta con sé una
replicazione incorporata perché è osservato sotto entrambe le
architetture, ma nessuna inferenza statistica formale è possibile a
questo $n$, e nessuna è rivendicata.

\paragraph{Rumore di valutazione, reward e mappe.}
In valutazione, 100 episodi per scenario significano 300 record
veicolo: a tassi di fascia media il solo errore standard binomiale è
$\approx$3\,pp per scenario, quindi le differenze di pochi punti a
livello di singolo scenario non sono individualmente significative e si
interpretano solo i pattern tra scenari. La funzione di reward, inoltre,
è stata calibrata attraverso iterazioni eseguite per lo più su
configurazioni bloccate su easy, quindi lo shaping del trunk può essere
implicitamente sbilanciato verso i percorsi corti; entrambi i regimi
condividono però la reward identica, quindi il contrasto primario resta
internamente valido e l'avvertenza riguarda la sola prestazione
\emph{assoluta} sui percorsi hard. La generalizzazione, infine, è
misurata su esattamente un trasferimento, e le affermazioni sono
limitate a questa coppia di mappe. Va anche ricordato che la pipeline di
valutazione stessa ha fallito la validità una volta e che il difetto di
lunghezza dei percorsi ha ridefinito silenziosamente la distribuzione
dei task per un'intera era di esperimenti (\cref{sec:validita}): in una
piattaforma custom, gli effetti apparenti devono prima sopravvivere alla
domanda ``lo strumento sta misurando ciò che dichiara?''.

\paragraph{Limiti di scope.}
Il design fissa infine diverse dimensioni, e i risultati non vanno letti
oltre di esse (\cref{sec:campagna}). L'architettura degli \emph{actor} è
fissa, quindi l'asse GNN parla della sola stima del valore sotto CTDE; la
cornice a mondo condiviso non è ablata, quindi il suo contributo causale
è caratterizzato e non isolato; le osservazioni sono privilegiate e non
percettive, quindi nulla segue per pipeline basate su camera o lidar; e
l'algoritmo e la versione della piattaforma sono unici, quindi altri
algoritmi possono interagire diversamente con l'ordinamento.
